\chapter{Complexity — Encoding}
%%% generated by the blueprint pipeline from TCSlib.Complexity.CircuitComplexity.Encoding
\section{Overview}

Declarations of the Lean module \texttt{TCSlib.Complexity.CircuitComplexity.Encoding}.

\section{Declarations}

\begin{definition}[Encoding a Boolean circuit tree]
\lean{BoolCircuit.encodeCircuit}
\label{BoolCircuit.encodeCircuit}
\leanok
\uses{BoolCircuit.TreeCircuit, BoolCircuit.encodeNat}
Encode a literal with index $i$ and polarity $s$ as a false bit, followed by $s$,
followed by $i$ true bits and then a false bit. Encode an AND or OR node by a true bit,
its Boolean gate label, the encodings of its children each preceded by a true bit, and a
final false bit.
\end{definition}

\begin{definition}[Encoding a list of child circuits]
\lean{BoolCircuit.encodeChildren}
\label{BoolCircuit.encodeChildren}
\leanok
\uses{BoolCircuit.TreeCircuit, BoolCircuit.encodeCircuit}
For a list of child circuits, encode each child by its circuit encoding preceded by a
true bit, concatenate these blocks in order, and append a false bit.
\end{definition}

\begin{theorem}[Encoding an internal circuit node]
\lean{BoolCircuit.encodeCircuit_node}
\label{BoolCircuit.encodeCircuit_node}
\leanok
\uses{BoolCircuit.TreeCircuit, BoolCircuit.encodeChildren, BoolCircuit.encodeCircuit}
\difficulty{1}
For any natural number $n$, any Boolean gate choice $b$, and any list $C$ of circuit
trees on $n$ variables, the encoding of the node with gate choice $b$ and children $C$
is the bit $1$, followed by $b$, followed by the encoding of the child list $C$.
\end{theorem}

\begin{theorem}[Encoding an empty child list]
\lean{BoolCircuit.encodeChildren_nil}
\label{BoolCircuit.encodeChildren_nil}
\leanok
\uses{BoolCircuit.TreeCircuit, BoolCircuit.encodeChildren}
\difficulty{1}
For any natural number $n$, the encoding of the empty list of circuit trees on $n$
variables is the one-bit list $[0]$.
\end{theorem}

\begin{theorem}[Encoding a nonempty child list]
\lean{BoolCircuit.encodeChildren_cons}
\label{BoolCircuit.encodeChildren_cons}
\leanok
\uses{BoolCircuit.TreeCircuit, BoolCircuit.encodeChildren, BoolCircuit.encodeCircuit}
\difficulty{1}
For any natural number $n$, any Boolean circuit tree $d$ on $n$ variables, and any list
$D$ of such trees, the encoding of the child list beginning with $d$ is the bit $1$,
followed by the encoding of $d$, followed by the encoding of $D$.
\end{theorem}

\begin{definition}[Reading a Boolean circuit]
\lean{BoolCircuit.readCircuit}
\label{BoolCircuit.readCircuit}
\leanok
\uses{BoolCircuit.TreeCircuit, BoolCircuit.decodeNat}
Given a natural number $n$, a fuel bound, and a Boolean list, the circuit reader returns
either no result or a circuit on $n$ variables together with the unconsumed suffix of
the list. With zero fuel it returns no result; with positive fuel it reads a literal
from a false bit, a polarity bit, and a unary index encoding when the decoded index is
less than $n$, or reads a node from a true bit, a gate label, and a valid encoding of
its children.
\end{definition}

\begin{definition}[Reading a list of circuit children]
\lean{BoolCircuit.readChildren}
\label{BoolCircuit.readChildren}
\leanok
\uses{BoolCircuit.TreeCircuit}
For a variable bound $n$, a fuel bound, and a Boolean list, the child-list reader
returns either no result or a list of Boolean circuit trees together with the unused
suffix. It fails with zero fuel or on malformed input; a leading false ends the list,
while a leading true signals another circuit followed recursively by the remaining
children.
\end{definition}

\begin{theorem}[Decoding a list of child circuits]
\lean{BoolCircuit.readChildren_encodeChildren}
\label{BoolCircuit.readChildren_encodeChildren}
\leanok
\uses{BoolCircuit.TreeCircuit, BoolCircuit.encodeChildren, BoolCircuit.encodeChildren_cons, BoolCircuit.encodeChildren_nil, BoolCircuit.encodeCircuit, BoolCircuit.readChildren, BoolCircuit.readCircuit}
\difficulty{5}
Let $n$ be a natural number and let $ds$ be a list of Boolean circuit trees on $n$
variables. Suppose each tree $d$ in $ds$ can be read back from its encoding whenever the
available fuel is at least the encoding length, leaving any appended remainder
unchanged. Then, for every fuel and remainder, if the encoding of the whole child list
has length at most the fuel, reading that encoding followed by the remainder returns
exactly $ds$ and the unchanged remainder.
\end{theorem}

\begin{theorem}[Circuit encoding and decoding]
\lean{BoolCircuit.readCircuit_encodeCircuit}
\label{BoolCircuit.readCircuit_encodeCircuit}
\leanok
\uses{BoolCircuit.TreeCircuit, BoolCircuit.TreeCircuit.ind, BoolCircuit.decodeNat_encodeNat, BoolCircuit.encodeCircuit, BoolCircuit.encodeCircuit_node, BoolCircuit.readChildren_encodeChildren, BoolCircuit.readCircuit}
\difficulty{4}
For every Boolean circuit $C$ on $n$ variables, every fuel bound $f$, and every Boolean
list $r$, if the length of the encoding of $C$ is at most $f$, then reading the encoded
circuit followed by $r$ with fuel $f$ returns exactly $C$ and leaves $r$ unread.
\end{theorem}

\begin{definition}[Encoding a circuit with its variable bound]
\lean{BoolCircuit.encodeSigma}
\label{BoolCircuit.encodeSigma}
\leanok
\uses{BoolCircuit.TreeCircuit, BoolCircuit.encodeCircuit, BoolCircuit.encodeNat}
For each natural number $n$, encode a Boolean circuit tree on $n$ variables by
concatenating the unary Boolean encoding of $n$ with the encoding of the circuit tree.
\end{definition}

\begin{definition}[Decoding a circuit with its variable bound]
\lean{BoolCircuit.decodeSigma}
\label{BoolCircuit.decodeSigma}
\leanok
\uses{BoolCircuit.TreeCircuit, BoolCircuit.decodeNat, BoolCircuit.encodeSigma, BoolCircuit.readCircuit}
Given a Boolean list, first decode its initial unary natural number $n$, then read a
circuit tree on $n$ variables from the remaining list. Return the resulting pair only if
the reader consumes the entire remainder and re-encoding the pair gives exactly the
original Boolean list; otherwise return no result.
\end{definition}

\begin{theorem}[Recovery of encoded circuit pairs]
\lean{BoolCircuit.decodeSigma_encodeSigma}
\label{BoolCircuit.decodeSigma_encodeSigma}
\leanok
\uses{BoolCircuit.TreeCircuit, BoolCircuit.decodeNat_encodeNat, BoolCircuit.decodeSigma, BoolCircuit.encodeCircuit, BoolCircuit.encodeSigma, BoolCircuit.readCircuit, BoolCircuit.readCircuit_encodeCircuit}
\difficulty{3}
For every natural number $n$ and every Boolean circuit tree $C$ on $n$ variables,
decoding the encoding of the pair $(n,C)$ returns that same pair.
\end{theorem}

\begin{theorem}[Encoding after successful decoding]
\lean{BoolCircuit.encodeSigma_of_decodeSigma}
\label{BoolCircuit.encodeSigma_of_decodeSigma}
\leanok
\uses{BoolCircuit.TreeCircuit, BoolCircuit.decodeSigma, BoolCircuit.encodeSigma}
\difficulty{4}
For any Boolean list $bs$ and any pair consisting of a natural number $n$ and a tree
circuit on $n$ variables, if decoding $bs$ succeeds and returns that pair, then encoding
the pair gives $bs$.
\end{theorem}

\begin{definition}[Boolean encoding of bounded circuit trees]
\lean{BoolCircuit.circuitEncoding}
\label{BoolCircuit.circuitEncoding}
\leanok
\uses{BoolCircuit.TreeCircuit, BoolCircuit.decodeSigma, BoolCircuit.decodeSigma_encodeSigma, BoolCircuit.encodeSigma}
Equip the type of pairs $(n,C)$, where $n$ is a natural number and $C$ is a Boolean
circuit tree on $n$ variables, with a finite encoding over the Boolean alphabet. Its
encoder is the concatenated unary encoding of $n$ and the circuit-tree encoding, and its
decoder is the partial decoder that accepts only fully consumed, canonically encoded
lists.
\end{definition}

\begin{theorem}[Circuit encoding map]
\lean{BoolCircuit.circuitEncoding_encode}
\label{BoolCircuit.circuitEncoding_encode}
\leanok
\uses{BoolCircuit.circuitEncoding, BoolCircuit.encodeSigma}
\difficulty{1}
The encoding map for pairs consisting of a natural-number variable bound and a Boolean
circuit tree is the map that concatenates the unary encoding of the bound with the
encoding of the circuit tree.
\end{theorem}

\begin{theorem}[Circuit encoding and decoding]
\lean{BoolCircuit.circuitEncoding_decode}
\label{BoolCircuit.circuitEncoding_decode}
\leanok
\uses{BoolCircuit.circuitEncoding, BoolCircuit.decodeSigma}
\difficulty{1}
For every finite variable bound, the decoder of the circuit encoding agrees with the
decoder that first reads the variable bound and then reads a circuit, accepting only
when re-encoding the result recovers the input bit string.
\end{theorem}

\begin{definition}[Language of satisfiable circuit encodings]
\lean{BoolCircuit.cktSatLang}
\label{BoolCircuit.cktSatLang}
\leanok
\uses{BoolCircuit.CktSat, BoolCircuit.decodeSigma}
The language consists of Boolean words that decode to some circuit belonging to the
class of satisfiable circuits.
\end{definition}

\begin{theorem}[Circuit Encoding and Satisfiability]
\lean{BoolCircuit.mem_cktSatLang_iff}
\label{BoolCircuit.mem_cktSatLang_iff}
\leanok
\uses{BoolCircuit.CktSat, BoolCircuit.TreeCircuit, BoolCircuit.cktSatLang, BoolCircuit.decodeSigma_encodeSigma, BoolCircuit.encodeSigma}
\difficulty{3}
For a Boolean circuit tree on $n$ variables, its encoding as a bit string belongs to the
language of encodings of satisfiable circuit trees if and only if the circuit tree is
satisfiable.
\end{theorem}

\begin{theorem}[Characterization of the Circuit Satisfiability Language]
\lean{BoolCircuit.mem_cktSatLang_iff_exists}
\label{BoolCircuit.mem_cktSatLang_iff_exists}
\leanok
\uses{BoolCircuit.CktSat, BoolCircuit.cktSatLang, BoolCircuit.encodeSigma, BoolCircuit.encodeSigma_of_decodeSigma, BoolCircuit.mem_cktSatLang_iff}
\difficulty{3}
A bit string belongs to the language of encodings of satisfiable Boolean circuit trees
if and only if it is the encoding of some satisfiable Boolean circuit tree.
\end{theorem}

\begin{theorem}[Circuit size and encoding length]
\lean{BoolCircuit.size_le_length_encodeCircuit}
\label{BoolCircuit.size_le_length_encodeCircuit}
\leanok
\uses{BoolCircuit.TreeCircuit, BoolCircuit.TreeCircuit.ind, BoolCircuit.TreeCircuit.size, BoolCircuit.encodeChildren, BoolCircuit.encodeChildren_cons, BoolCircuit.encodeChildren_nil, BoolCircuit.encodeCircuit, BoolCircuit.encodeCircuit_node, BoolCircuit.encodeNat}
\difficulty{5}
For any natural number $n$ and any Boolean circuit tree on $n$ variables, its size is at
most the length of its Boolean-list encoding.
\end{theorem}

\begin{theorem}[Circuit size and paired encoding length]
\lean{BoolCircuit.size_le_length_encodeSigma}
\label{BoolCircuit.size_le_length_encodeSigma}
\leanok
\uses{BoolCircuit.TreeCircuit, BoolCircuit.TreeCircuit.size, BoolCircuit.encodeNat, BoolCircuit.encodeSigma, BoolCircuit.size_le_length_encodeCircuit}
\difficulty{3}
For any natural number $n$ and Boolean circuit tree $C$ on $n$ variables, the size of
$C$ is at most the length of the encoding formed by concatenating the unary encoding of
$n$ with the Boolean-list encoding of $C$.
\end{theorem}

\begin{theorem}[Total length of gate clauses]
\lean{BoolCircuit.gateClauses_flatten_length}
\label{BoolCircuit.gateClauses_flatten_length}
\leanok
\uses{BoolCircuit.CktVar, BoolCircuit.TreeCircuit, BoolCircuit.gateClauses, SATTo3SAT.Literal}
\difficulty{4}
Let $n$ be a natural number, let $b$ be a Boolean value, and let $cs$ be a finite list
of Boolean circuit trees on $n$ variables. The total number of literal occurrences in
the clauses associated with $b$ and $cs$ is $3|cs|+1$.
\end{theorem}

\begin{theorem}[Tseitin Flattening Length Bound]
\lean{BoolCircuit.tseitin_flatten_length_le}
\label{BoolCircuit.tseitin_flatten_length_le}
\leanok
\uses{BoolCircuit.TreeCircuit, BoolCircuit.TreeCircuit.ind, BoolCircuit.TreeCircuit.size, BoolCircuit.TreeCircuit.tseitin, BoolCircuit.gateClauses_flatten_length}
\difficulty{5}
For every Boolean circuit tree $C$ on $n$ variables, the length of the flattened Tseitin
formula associated with $C$, plus $3$, is at most $7$ times the size of $C$.
\end{theorem}

\begin{theorem}[Length bound for the flattened CNF]
\lean{BoolCircuit.toCNF_flatten_length_le}
\label{BoolCircuit.toCNF_flatten_length_le}
\leanok
\uses{BoolCircuit.TreeCircuit, BoolCircuit.TreeCircuit.size, BoolCircuit.TreeCircuit.toCNF, BoolCircuit.tseitin_flatten_length_le}
\difficulty{3}
For every Boolean circuit tree $C$, the length of the flattened list of clauses in its
conjunctive normal form is at most $7\,\mathrm{size}(C)$.
\end{theorem}

\begin{theorem}[Size bound for the 3-CNF encoding]
\lean{BoolCircuit.length_to3SAT_toCNF_le}
\label{BoolCircuit.length_to3SAT_toCNF_le}
\leanok
\uses{BoolCircuit.TreeCircuit, BoolCircuit.TreeCircuit.toCNF, BoolCircuit.TreeCircuit.toCNF_length_le, BoolCircuit.encodeSigma, BoolCircuit.size_le_length_encodeSigma, BoolCircuit.toCNF_flatten_length_le, SATTo3SAT.length_to3SAT_le, SATTo3SAT.to3SAT}
\difficulty{3}
Let $n\in\mathbb{N}$ and let $C$ be a Boolean circuit tree on $n$ variables. The number
of 3-clauses in the 3-SAT encoding of the circuit’s Tseitin CNF is at most $13$ times
the length of the encoding of the pair $(n,C)$.
\end{theorem}
